\begin{proof}[Proof of \cref{obj:lem:factorial-derivatives}]
\begin{enumerate}
\item We first treat \(j=1\), recording identities that will also be used for higher derivatives. Define, for every \(k\ge0\),
\[
R_k(x)=P_k^{\mathrm{Leg}}(2x-1),\qquad
\omega_k=2k+1,\qquad
c_k=\omega_k\int_0^1 p(x)R_k(x)\,dx,
\qquad
\mathcal E_p=\int_0^1 p(x)^2\,dx.
\]
Changing variables in \cref{obj:lem:classical-legendre-facts} gives
\[
\int_0^1 R_k(x)R_l(x)\,dx
=
\begin{cases}
1/\omega_k,&k=l,\\
0,&k\ne l.
\end{cases}
\]
The polynomials \((-1)^kR_k\), \(0\le k\le m\), have respective degrees \(0,\ldots,m\), so they span the polynomials of degree at most \(m\). The same is therefore true of the \(R_k\). Integrating a basis expansion against \(R_k\) identifies its coefficient as \(c_k\). Consequently,
\[
p(x)=\sum_{k=0}^m c_kR_k(x),
\qquad
\mathcal E_p=\sum_{k=0}^m\frac{c_k^2}{\omega_k}\ge0.
\]
The second identity follows by expanding the square and applying orthogonality.
% lean: CausalSmith.Stat.RdTruesideNoiseFrontier.polynomial_shiftedLegendre_parseval

\item Differentiating the expansion and applying weighted Cauchy--Schwarz yields, for every real \(x\),
\[
p'(x)=\sum_{k=0}^m c_kR_k'(x),
\qquad
p'(x)^2
\le
\left(\sum_{k=0}^m\frac{c_k^2}{\omega_k}\right)
\left(\sum_{k=0}^m\omega_kR_k'(x)^2\right)
=
\mathcal E_p\sum_{k=0}^m\omega_kR_k'(x)^2.
\]
To integrate this bound, first observe that
\[
\sum_{i=0}^{k-1}\omega_i=k^2\qquad(k\ge0);
\]
this follows by induction, since adding \(\omega_k=2k+1\) changes \(k^2\) to \((k+1)^2\). Define the parity-restricted sums
\[
\pi_k=\sum_{\substack{0\le i<k\\k-i\ \mathrm{odd}}}\omega_i
\qquad(k\ge0).
\]
The parity classes alternate when \(k\) increases, and the new index \(i=k\) belongs to the sum for \(k+1\). Thus
\[
\pi_0=0,\qquad
\pi_{k+1}+\pi_k=(k+1)^2,
\qquad
2\pi_k=k(k+1),
\]
where the last identity follows by induction from the preceding recurrence.

For \(k\ge1\), \cref{obj:lem:shifted-legendre-derivative} supplies the derivative expansion. Orthogonality then gives
\[
\int_0^1 R_k'(x)^2\,dx
=
4\sum_{\substack{0\le i<k\\k-i\ \mathrm{odd}}}\omega_i
=
2k(k+1).
\]
For \(k=0\), the same identity holds because \(R_0(x)=1\) and \(R_0'(x)=0\). Integrating the pointwise bound therefore gives
\[
\int_0^1 p'(x)^2\,dx
\le
\mathcal E_p\sum_{k=0}^m\omega_k\,2k(k+1).
\]
Since \(0\le k\le m\) implies \(k(k+1)\le(m+1)^2\),
\[
\sum_{k=0}^m\omega_k\,2k(k+1)
\le
2(m+1)^2\sum_{k=0}^m\omega_k
=
2(m+1)^4
\le
4(m+1)^4.
\]
It follows that
\[
\int_0^1 p'(x)^2\,dx
\le
4(m+1)^4\mathcal E_p
\le
16(m+1)^4\mathcal E_p
=
\left(4(m+1)^2\sqrt{\mathcal E_p}\right)^2.
\]
Taking nonnegative square roots proves the required inequality when \(j=1\).
% lean: CausalSmith.Stat.RdTruesideNoiseFrontier.polynomial_first_derivative_norm_bound

\item Now suppose \(j\ne1\). Since \(j\ge1\), this means \(j\ge2\). Define the nonnegative constant
\[
\Lambda_{m,j}=\frac{4^j(m+1)^{2j}}{j!}.
\]
On the index set \(0,\ldots,m\), define the differentiation matrix
\[
(\mathsf J_m)_{kl}
=
\begin{cases}
2\omega_k,&k<l\ \text{and }l-k\text{ is odd},\\
0,&\text{otherwise}.
\end{cases}
\]
Its powers are defined by
\[
(\mathsf J_m^0)_{kl}
=
\begin{cases}
1,&k=l,\\
0,&k\ne l,
\end{cases}
\qquad
(\mathsf J_m^{r+1})_{kl}
=
\sum_{i=0}^m(\mathsf J_m^r)_{ki}(\mathsf J_m)_{il}
\quad(r\ge0).
\]
Also define the derivative coefficients, for \(r\ge0\) and \(0\le k\le m\), by
\[
c_k^{[r]}=\sum_{l=0}^m(\mathsf J_m^r)_{kl}c_l.
\]
The derivative expansion in \cref{obj:lem:shifted-legendre-derivative}, together with \(R_0'=0\), gives
\[
R_l'(x)=\sum_{k=0}^m(\mathsf J_m)_{kl}R_k(x)
\qquad(0\le l\le m).
\]
Differentiating finite sums applies this coefficient matrix once; induction on \(r\) therefore gives
\[
p^{(r)}(x)=\sum_{k=0}^m c_k^{[r]}R_k(x).
\]
Orthogonality now identifies both square integrals:
\[
\int_0^1 p(x)^2\,dx
=
\sum_{k=0}^m\frac{c_k^2}{\omega_k},
\qquad
\int_0^1 p^{(j)}(x)^2\,dx
=
\sum_{k=0}^m\frac{(c_k^{[j]})^2}{\omega_k}.
\]
% lean: CausalSmith.Stat.RdTruesideNoiseFrontier.polynomial_iterated_derivative_parseval

\item We next bound the entries of the matrix powers. The prefix estimate needed below is
\[
(r+1)\sum_{i=0}^{k-1}\omega_i(i^2)^r
\le
(k^2)^{r+1}
\qquad(k,r\ge0).
\]
To establish it, use the elementary inequality
\[
\xi^{r+1}+(r+1)\zeta\xi^r
\le
(\xi+\zeta)^{r+1}
\qquad(\xi,\zeta\ge0,\ r\ge0).
\]
For \(r=0\) this is equality. Its inductive step follows from
\[
\xi^{r+2}+(r+2)\zeta\xi^{r+1}
\le
\bigl(\xi^{r+1}+(r+1)\zeta\xi^r\bigr)(\xi+\zeta)
\le
(\xi+\zeta)^{r+2},
\]
because the difference in the first inequality is
\((r+1)\zeta^2\xi^r\ge0\).
For fixed \(r\), the prefix estimate follows by induction on \(k\). The empty sum at \(k=0\) is zero, and the inductive step is
\[
\begin{aligned}
(r+1)\sum_{i=0}^{k}\omega_i(i^2)^r
&\le
(k^2)^{r+1}+(r+1)(2k+1)(k^2)^r\\
&\le
\bigl(k^2+2k+1\bigr)^{r+1}
=
\bigl((k+1)^2\bigr)^{r+1}.
\end{aligned}
\]
Here, as usual, the zeroth power is \(1\), including at zero.
% lean: CausalSmith.Stat.RdTruesideNoiseFrontier.shiftedLegendre_prefix_power_bound

\item Every entry of every matrix power is nonnegative, by its defining recurrence. We claim that
\[
0\le(\mathsf J_m^r)_{kl}
\le
\frac{2^r\omega_k(l^2)^{r-1}}{(r-1)!}
\qquad(r\ge1,\ 0\le k,l\le m).
\]
For \(r=1\), this follows directly from
\(0\le(\mathsf J_m)_{kl}\le2\omega_k\).
Suppose the bound holds for \(r\). Entries \((\mathsf J_m)_{il}\) vanish unless \(i<l\), and are bounded by \(2\omega_i\). The recurrence and the prefix estimate therefore give
\[
\begin{aligned}
(\mathsf J_m^{r+1})_{kl}
&\le
\sum_{i=0}^{l-1}
\frac{2^r\omega_k(i^2)^{r-1}}{(r-1)!}\,2\omega_i\\
&=
\frac{2^{r+1}\omega_k}{(r-1)!}
\sum_{i=0}^{l-1}\omega_i(i^2)^{r-1}\\
&\le
\frac{2^{r+1}\omega_k(l^2)^r}{r(r-1)!}
=
\frac{2^{r+1}\omega_k(l^2)^r}{r!}.
\end{aligned}
\]
This also covers \(l=0\), for which the sum is empty. Define
\[
\Gamma_{m,j}
=
\frac{2^j(m+1)^{2(j-1)}}{(j-1)!}\ge0.
\]
Since \(0\le l\le m\), the entry estimate implies
\[
0\le(\mathsf J_m^j)_{kl}\le\Gamma_{m,j}\omega_k,
\qquad
\bigl((\mathsf J_m^j)_{kl}\bigr)^2
\le
\Gamma_{m,j}^2\omega_k^2.
\]
% lean: CausalSmith.Stat.RdTruesideNoiseFrontier.legendreDerivativePower_entry_bound

\item Weighted Cauchy--Schwarz transfers this entry bound to the coefficients. For every \(0\le k\le m\),
\[
\begin{aligned}
(c_k^{[j]})^2
&\le
\left(\sum_{l=0}^m
\bigl((\mathsf J_m^j)_{kl}\bigr)^2\omega_l\right)
\left(\sum_{l=0}^m\frac{c_l^2}{\omega_l}\right)\\
&\le
\Gamma_{m,j}^2\omega_k^2
\left(\sum_{l=0}^m\omega_l\right)
\left(\sum_{l=0}^m\frac{c_l^2}{\omega_l}\right).
\end{aligned}
\]
Dividing by the positive weight \(\omega_k\), summing over \(k\), and using the total-weight identity gives
\[
\sum_{k=0}^m\frac{(c_k^{[j]})^2}{\omega_k}
\le
\Gamma_{m,j}^2
\left(\sum_{k=0}^m\omega_k\right)^2
\left(\sum_{l=0}^m\frac{c_l^2}{\omega_l}\right)
=
\Gamma_{m,j}^2(m+1)^4
\sum_{l=0}^m\frac{c_l^2}{\omega_l}.
\]
% lean: CausalSmith.Stat.RdTruesideNoiseFrontier.weighted_derivative_matrix_action_bound

\item For every integer \(r\ge0\), induction gives \(r\le2^r\): the base case is immediate, and
\[
r+1\le2^r+1\le2^{r+1}.
\]
Using \(j!=j(j-1)!\) and \(4^j=2^j2^j\), we obtain
\[
\frac{2^j}{(j-1)!}
=
\frac{j\,2^j}{j!}
\le
\frac{4^j}{j!}.
\]
Consequently,
\[
\Gamma_{m,j}(m+1)^2
=
\frac{2^j}{(j-1)!}(m+1)^{2j}
\le
\Lambda_{m,j}.
\]
Both sides are nonnegative, so squaring and substituting into the coefficient estimate yields
\[
\sum_{k=0}^m\frac{(c_k^{[j]})^2}{\omega_k}
\le
\Lambda_{m,j}^2
\sum_{k=0}^m\frac{c_k^2}{\omega_k}.
\]
% lean: CausalSmith.Stat.RdTruesideNoiseFrontier.legendreDerivativeAction_factorial_bound

\item Finally, the two square-integral identities give
\[
\int_0^1 p^{(j)}(x)^2\,dx
\le
\Lambda_{m,j}^2\mathcal E_p
=
\left(\Lambda_{m,j}\sqrt{\mathcal E_p}\right)^2.
\]
Since \(\mathcal E_p\ge0\) and \(\Lambda_{m,j}\ge0\), taking nonnegative square roots proves
\[
\left(\int_0^1 p^{(j)}(x)^2\,dx\right)^{1/2}
\le
\frac{4^j(m+1)^{2j}}{j!}
\left(\int_0^1 p(x)^2\,dx\right)^{1/2}.
\]
The argument includes \(\mathcal E_p=0\) and the endpoint order \(j=m\).
% lean: CausalSmith.Stat.RdTruesideNoiseFrontier.factorial_derivatives
\end{enumerate}
\end{proof}
